\begin{frame}[fragile]\frametitle{}
\begin{center}
{\Large LLMs, SLLMs \& Agents}

{\small How they work, which to use, and when}
\end{center}
\end{frame}

%%%%%%%%%%%%%%%%%%%%%%%%%%%%%%%%%%%%%%%%%%%%%%%%%%%%%%%%%%%
\begin{frame}[fragile]\frametitle{What Is a Large Language Model (LLM)?}
\textbf{How it works (college-level intuition):}
\begin{itemize}
\item An LLM is trained to predict: \textit{``Given these words, what word comes next?''}
\item Mathematically: given token sequence $w_1, w_2, \ldots, w_n$, compute $P(w_{n+1} \mid w_1 \ldots w_n)$
\item \textbf{What this means in plain English:} The model reads billions of web pages and books, and builds an internal ``map'' of which words and ideas go together. When you ask it a question, it generates the most likely continuation of your text.
\item During training on trillions of tokens, the model learns grammar, facts, reasoning patterns, and code, all from next-word prediction alone.
\end{itemize}
\textbf{Size matters:} ``Large'' = billions of parameters (weights). GPT-4 $\approx$ 1 trillion parameters.
\end{frame}

%%%%%%%%%%%%%%%%%%%%%%%%%%%%%%%%%%%%%%%%%%%%%%%%%%%%%%%%%%%
\begin{frame}[fragile]\frametitle{How Transformers Work, The Intuition}
\textbf{The key innovation: Attention Mechanism}
\begin{itemize}
\item Before Transformers: RNNs read text word-by-word, forgetting early context
\item Transformers read the \textit{entire sequence at once} and compute which words to \textit{pay attention to}
\item \textbf{Formula:} $\text{Attention}(Q, K, V) = \text{softmax}\!\left(\frac{QK^T}{\sqrt{d_k}}\right)V$
\item \textbf{What this means:} For each word, calculate a ``relevance score'' with every other word, then mix their meanings proportionally. The word ``bank'' in ``river bank'' pays more attention to ``river'' than ``money.''
\end{itemize}
\textbf{Result:} Transformers handle long documents, code, and multi-turn conversations much better than older models.
\end{frame}

%%%%%%%%%%%%%%%%%%%%%%%%%%%%%%%%%%%%%%%%%%%%%%%%%%%%%%%%%%%
\begin{frame}[fragile]\frametitle{The Modern LLM Landscape}
\textbf{Proprietary (API-access, closed weights):}
\begin{itemize}
\item \textbf{OpenAI:} GPT-4o, o1, o3, best general reasoning
\item \textbf{Anthropic:} Claude 3.5 / 3.7 Sonnet, best for coding, long context
\item \textbf{Google:} Gemini 2.0 Flash, native multimodal, fast, cheap
\item \textbf{Microsoft:} Azure OpenAI, enterprise-grade, compliance
\end{itemize}

\textbf{Open Weights (self-hostable):}
\begin{itemize}
\item \textbf{Meta:} Llama 3.3 70B, strong open alternative
\item \textbf{Mistral:} Mistral Large 2, Codestral, great for code
\item \textbf{DeepSeek:} DeepSeek-V3, strong reasoning, low cost
\end{itemize}

\textbf{Key insight for leaders:} Models are commoditizing. Your competitive advantage is in \textit{how you integrate them}, not which model you pick.
\end{frame}

%%%%%%%%%%%%%%%%%%%%%%%%%%%%%%%%%%%%%%%%%%%%%%%%%%%%%%%%%%%
\begin{frame}[fragile]\frametitle{Small LLMs (SLLMs), When Smaller Is Better}
\textbf{What is an SLLM?} A compact model (1B--13B parameters) designed for specific tasks or on-device deployment.

\begin{itemize}
\item \textbf{Examples:} Phi-3 Mini, Gemma 2B, Llama 3.2 3B, Mistral 7B
\item \textbf{When to use SLLM over LLM:}
  \begin{itemize}
  \item Latency is critical (real-time, embedded systems)
  \item Data must not leave your network (on-premise requirement)
  \item Task is narrow and well-defined (e.g., classify support ticket type)
  \item Cost is a primary constraint (1000$\times$ cheaper per token)
  \end{itemize}
\item \textbf{When to use a large LLM:}
  \begin{itemize}
  \item Complex multi-step reasoning required
  \item Long context window needed (legal docs, large codebases)
  \item You need the best possible quality without cost sensitivity
  \end{itemize}
\end{itemize}
\textbf{Rule:} Start with the smallest model that meets your quality bar.
\end{frame}

%%%%%%%%%%%%%%%%%%%%%%%%%%%%%%%%%%%%%%%%%%%%%%%%%%%%%%%%%%%
\begin{frame}[fragile]\frametitle{What Is an AI Agent?}
\textbf{Chatbot vs.\ Agent, the critical distinction:}
\begin{itemize}
\item \textbf{Chatbot:} One-shot Q\&A. User asks $\rightarrow$ LLM replies. No memory, no tools, no action.
\item \textbf{Agent:} LLM given \textit{tools} (web search, code execution, APIs, databases) and a \textit{goal}. It \textbf{plans}, executes steps, observes results, and adapts, autonomously.
\end{itemize}

\textbf{Agent Loop (ReAct pattern):}
\begin{center}
Observe $\rightarrow$ Think $\rightarrow$ Act $\rightarrow$ Observe $\rightarrow$ \ldots $\rightarrow$ Goal Achieved
\end{center}

\textbf{Engineering examples:}
\begin{itemize}
\item ``Analyze this failing build, find root cause in logs, create a JIRA ticket.'', 3 tool calls, no human in loop
\item ``Review this PR for security issues, check OWASP top 10, comment inline.'', Agent with code analysis tool
\end{itemize}
\end{frame}

%%%%%%%%%%%%%%%%%%%%%%%%%%%%%%%%%%%%%%%%%%%%%%%%%%%%%%%%%%%
\begin{frame}[fragile]\frametitle{Agentic Workflows Across the SDLC}
\begin{itemize}
\item \textbf{Requirements \& Design:}
  Requirement refinement agent, takes user stories, asks clarifying questions, generates acceptance criteria, flags ambiguities
\item \textbf{Coding:}
  Coding agent, given a ticket, writes code, runs tests locally, raises PR (e.g., Devin, Claude Code)
\item \textbf{Test Generation:}
  Agent generates unit tests, integration tests, edge cases from function signatures or API specs
\item \textbf{Code Review:}
  PR review agent, checks style, logic bugs, security (OWASP), performance, adds inline comments
\item \textbf{Release Notes:}
  Agent reads merged PRs + commit messages $\rightarrow$ drafts structured release notes
\item \textbf{Ops \& Support:}
  RCA agent, reads logs, traces, metrics $\rightarrow$ identifies root cause $\rightarrow$ suggests fix
\end{itemize}
\end{frame}

%%%%%%%%%%%%%%%%%%%%%%%%%%%%%%%%%%%%%%%%%%%%%%%%%%%%%%%%%%%
\begin{frame}[fragile]\frametitle{Live Demo Concept: Requirement $\rightarrow$ Code $\rightarrow$ Test $\rightarrow$ Docs}
\textbf{The full SDLC in one agentic pipeline:}
\begin{enumerate}
\item \textbf{Input:} ``Build a REST endpoint that returns user profile by ID''
\item \textbf{Agent Step 1, Design:} LLM clarifies, generates OpenAPI spec
\item \textbf{Agent Step 2, Code:} LLM writes \texttt{FastAPI} endpoint with error handling
\item \textbf{Agent Step 3, Tests:} LLM writes \texttt{pytest} unit + integration tests
\item \textbf{Agent Step 4, Docs:} LLM writes README section + Swagger annotations
\item \textbf{Agent Step 5, Review:} LLM does security review (no SQL injection, input validation)
\end{enumerate}
\textbf{Time:} $\approx$ 3--5 minutes. Human oversight: review and approve.\\
\textbf{Key lesson:} The engineer's job shifts from \textit{writing} to \textit{reviewing and directing}.
\end{frame}

%%%%%%%%%%%%%%%%%%%%%%%%%%%%%%%%%%%%%%%%%%%%%%%%%%%%%%%%%%%
\begin{frame}[fragile]\frametitle{MCP, Model Context Protocol}
\textbf{What is MCP?} An open standard (Anthropic, 2024) that defines how LLMs connect to external tools and data sources.

\begin{itemize}
\item \textbf{Problem before MCP:} Every AI tool had custom integrations. Fragile, hard to maintain.
\item \textbf{MCP solution:} Standardized protocol, LLM $\leftrightarrow$ Tool (like USB for AI)
\item \textbf{MCP architecture:}
  \begin{itemize}
  \item \textbf{MCP Host:} The AI app (e.g., Claude, Cursor)
  \item \textbf{MCP Server:} A connector to a specific tool (GitHub, Jira, Postgres, filesystem)
  \item \textbf{MCP Client:} Bridge inside the host that speaks the protocol
  \end{itemize}
\item \textbf{Engineering relevance:} Build one MCP server for your internal tool, it works with any MCP-compatible AI agent
\end{itemize}

\textbf{Analogy:} Before: proprietary chargers for every phone. After: USB-C for all.
\end{frame}


%%%%%%%%%%%%%%%%%%%%%%%%%%%%%%%%%%%%%%%%%%%%%%%%%%%%%%%%%%%
